\documentclass[11pt,a4paper]{article}
\usepackage[utf8]{inputenc}
\usepackage[T1]{fontenc}
\usepackage[italian]{babel}
\usepackage{amsmath,amssymb,bm}
\usepackage{siunitx}
\usepackage{booktabs}
\usepackage{geometry}
\geometry{margin=25mm}
\usepackage[colorlinks=true,linkcolor=blue,citecolor=blue]{hyperref}

\sisetup{output-decimal-marker={,}}

\title{Modello barra--stringa ed analisi FEM 1D:\\
risposta risonante ad eccitazione armonica a \SI{2}{kHz}}
\author{D.~Bagnara}
\date{\today}

\newcommand{\dpar}[2]{\frac{\partial #1}{\partial #2}}

\begin{document}
\maketitle

\begin{abstract}
Si presenta il modello dinamico di una \emph{barra--stringa} (trave di
Eulero--Bernoulli soggetta a tensione assiale) e la sua discretizzazione agli
elementi finiti con elementi di Hermite. Si ricavano le frequenze naturali in
forma chiusa per il caso appoggiato--appoggiato, si costruisce la risposta
armonica forzata e si mostra che, con un'opportuna scelta della tensione, la
struttura entra in risonanza per un'eccitazione sinusoidale a
$f_0 = \SI{2}{kHz}$, con una banda di risonanza (\emph{spread}) determinata
dallo smorzamento modale. Il documento accompagna lo script MATLAB
\texttt{beam\_string\_fem.m}.
\end{abstract}

\section{Introduzione}
Il modello barra--stringa descrive elementi strutturali snelli in cui
coesistono due meccanismi di richiamo elastico trasversale: la rigidezza
flessionale $EI$ (comportamento da \emph{barra}, o trave) e la tensione
assiale $T$ (comportamento da \emph{stringa}, o corda tesa). \`E il modello
naturale per fili, trefoli e conduttori tesi eccitati da forze distribuite o
concentrate, ad esempio di origine elettromagnetica. L'interesse pratico \`e
duplice: (i) prevedere \emph{a quale frequenza} la struttura risuona, dato che
in risonanza l'ampiezza \`e limitata solo dallo smorzamento; (ii) quantificare
lo \emph{spread} di frequenza entro cui la risposta rimane amplificata, cio\`e
la banda a met\`a potenza attorno al picco.

\section{Modello continuo}
\subsection{Equazione del moto}
Sia $w(x,t)$ lo spostamento trasversale di una campata di lunghezza $L$, massa
per unit\`a di lunghezza $\rho A$, rigidezza flessionale $EI$ e tensione
assiale $T$ costante. Il bilancio di quantit\`a di moto trasversale fornisce
la PDE del quarto ordine
\begin{equation}
\rho A \,\dpar{^2 w}{t^2}
+ c\,\dpar{w}{t}
+ EI\,\dpar{^4 w}{x^4}
- T\,\dpar{^2 w}{x^2}
= f(x,t),
\label{eq:pde}
\end{equation}
dove $c$ \`e uno smorzamento viscoso equivalente ed $f(x,t)$ la forza
trasversale per unit\`a di lunghezza. I due casi limite sono immediati: per
$T \to 0$ si ritrova la trave di Eulero--Bernoulli, per $EI \to 0$ l'equazione
delle corde vibranti. Il rapporto adimensionale
\begin{equation}
\mu_n = \frac{EI\,\beta_n^2}{T}, \qquad \beta_n = \frac{n\pi}{L},
\end{equation}
misura, modo per modo, quanto il comportamento sia da trave ($\mu_n \gg 1$) o
da corda ($\mu_n \ll 1$): la rigidezza flessionale pesa sempre di pi\`u al
crescere dell'ordine modale, perch\'e entra con la quarta potenza del numero
d'onda.

\subsection{Frequenze naturali (appoggiato--appoggiato)}
Con condizioni al contorno $w(0)=w(L)=0$, $w''(0)=w''(L)=0$, le autofunzioni
sono $\phi_n(x)=\sin(\beta_n x)$ e la relazione di dispersione fornisce
\begin{equation}
\omega_n^2 \;=\; \beta_n^2\,\frac{T}{\rho A}
\;+\; \beta_n^4\,\frac{EI}{\rho A},
\qquad
f_n = \frac{\omega_n}{2\pi}.
\label{eq:disp}
\end{equation}
Il primo termine \`e il contributo di corda (lineare in $\beta_n$ sulla
frequenza), il secondo quello flessionale (quadratico in $\beta_n$): lo
spettro di una barra--stringa \emph{non \`e armonico}, i modi alti sono
progressivamente ``stirati'' verso l'alto rispetto ai multipli interi del
fondamentale.

\subsection{Sintonizzazione della risonanza a \texorpdfstring{$f_0$}{f0}}
Fissata la geometria, la \eqref{eq:disp} si inverte rispetto a $T$: imposto
che il modo $n^\ast$ cada esattamente su $\omega_0 = 2\pi f_0$,
\begin{equation}
T^\ast = \frac{\rho A\,\omega_0^2 - EI\,\beta_{n^\ast}^4}{\beta_{n^\ast}^2},
\qquad \beta_{n^\ast} = \frac{n^\ast\pi}{L},
\label{eq:tuning}
\end{equation}
con il vincolo di ammissibilit\`a $T^\ast > 0$ (e $\sigma = T^\ast/A$ entro il
limite del materiale). Nello script il modo $n^\ast$ \`e scelto come quello
pi\`u vicino a $f_0$ per la tensione di primo tentativo.

\section{Formulazione debole e discretizzazione FEM}
\subsection{Forma debole}
Moltiplicando la \eqref{eq:pde} per una funzione test $v \in H^2_0(0,L)$ ed
integrando per parti due volte il termine flessionale ed una volta quello di
tensione si ottiene: trovare $w(\cdot,t)$ tale che
\begin{equation}
\int_0^L \rho A\, \ddot{w}\, v\,\mathrm{d}x
+ \int_0^L c\, \dot{w}\, v\,\mathrm{d}x
+ \int_0^L EI\, w''\, v''\,\mathrm{d}x
+ \int_0^L T\, w'\, v'\,\mathrm{d}x
= \int_0^L f\, v\,\mathrm{d}x
\label{eq:weak}
\end{equation}
per ogni $v$ ammissibile. La presenza di $w''$ richiede continuit\`a $C^1$
delle funzioni di forma: la scelta canonica sono i polinomi cubici di Hermite.

\subsection{Elemento di Hermite}
Ogni elemento di lunghezza $h$ ha due nodi e quattro gradi di libert\`a
$\bm{q}_e = [\,w_1\;\theta_1\;w_2\;\theta_2\,]^{\mathsf T}$ con
$\theta_i = w'(x_i)$. In coordinate naturali $\xi = (x-x_1)/h \in [0,1]$,
\begin{equation}
\begin{aligned}
N_1 &= 1 - 3\xi^2 + 2\xi^3, &\quad N_2 &= h(\xi - 2\xi^2 + \xi^3),\\
N_3 &= 3\xi^2 - 2\xi^3,     &\quad N_4 &= h(-\xi^2 + \xi^3).
\end{aligned}
\end{equation}
Sostituendo $w = \bm{N}\bm{q}_e$ nella \eqref{eq:weak} si ottengono le tre
matrici elementari. Rigidezza flessionale:
\begin{equation}
\bm{K}_b^{(e)} = \frac{EI}{h^3}
\begin{bmatrix}
12 & 6h & -12 & 6h\\
6h & 4h^2 & -6h & 2h^2\\
-12 & -6h & 12 & -6h\\
6h & 2h^2 & -6h & 4h^2
\end{bmatrix}.
\end{equation}
Rigidezza \emph{geometrica} indotta dalla tensione (dal termine $T\,w'v'$):
\begin{equation}
\bm{K}_g^{(e)} = \frac{T}{30h}
\begin{bmatrix}
36 & 3h & -36 & 3h\\
3h & 4h^2 & -3h & -h^2\\
-36 & -3h & 36 & -3h\\
3h & -h^2 & -3h & 4h^2
\end{bmatrix}.
\end{equation}
Massa consistente:
\begin{equation}
\bm{M}^{(e)} = \frac{\rho A h}{420}
\begin{bmatrix}
156 & 22h & 54 & -13h\\
22h & 4h^2 & 13h & -3h^2\\
54 & 13h & 156 & -22h\\
-13h & -3h^2 & -22h & 4h^2
\end{bmatrix}.
\end{equation}
L'assemblaggio standard produce $\bm{M}$, $\bm{K} = \bm{K}_b + \bm{K}_g$
globali; le condizioni appoggiato--appoggiato si impongono eliminando i gradi
di libert\`a $w(0)$ e $w(L)$ (le rotazioni d'estremit\`a restano libere). Un
carico puntuale $F_0$ in $x_F$ si proietta sui gradi di libert\`a
dell'elemento contenente $x_F$ tramite le stesse funzioni di forma,
$\bm{F} = F_0\,\bm{N}^{\mathsf T}(\xi_F)$: \`e il carico
\emph{energeticamente consistente}.

\subsection{Smorzamento di Rayleigh}
Lo smorzamento \`e modellato come
$\bm{C} = \alpha\bm{M} + \beta\bm{K}$, con $\alpha,\beta$ calibrati imponendo
lo smorzamento modale $\zeta$ a due pulsazioni $\omega_1,\omega_2$ che
delimitano la banda d'interesse:
\begin{equation}
\alpha = \frac{2\zeta\,\omega_1\omega_2}{\omega_1+\omega_2},
\qquad
\beta = \frac{2\zeta}{\omega_1+\omega_2}.
\end{equation}
Nello script si usa $\omega_1 = \omega_0$, $\omega_2 = 2\omega_0$ con
$\zeta = 0{,}5\,\%$, valore tipico per fili metallici tesi.

\section{Analisi dinamica}
\subsection{Analisi modale}
Il problema agli autovalori generalizzato
$\bm{K}\bm{\phi}_n = \omega_n^2\,\bm{M}\bm{\phi}_n$ fornisce frequenze e forme
modali FEM, da confrontare con la \eqref{eq:disp} per la verifica di
convergenza. Con elementi di Hermite l'errore in frequenza converge come
$O(h^4)$: con $N_e = 200$ elementi su $L = \SI{0.8}{m}$ (oltre 20 elementi per
lunghezza d'onda del modo 9) l'errore sul modo target risulta
$\sim 2\cdot10^{-6}\,\%$.

\subsection{Risposta armonica forzata (FRF)}
Per $f(x,t) = F_0\,\delta(x-x_F)\sin(\omega t)$ la risposta a regime \`e
$\bm{w}(t) = \Im\{\bm{W}(\omega)e^{\mathrm{j}\omega t}\}$ con
\begin{equation}
\bm{W}(\omega) =
\bigl(\bm{K} + \mathrm{j}\omega\bm{C} - \omega^2\bm{M}\bigr)^{-1}\bm{F}.
\label{eq:frf}
\end{equation}
Espandendo sulla base modale ($\bm{\Phi}^{\mathsf T}\bm{M}\bm{\Phi}=\bm{I}$),
\begin{equation}
W(x_o,\omega) = \sum_n
\frac{\phi_n(x_o)\,\phi_n(x_F)\,F_0}
{\omega_n^2 - \omega^2 + 2\mathrm{j}\zeta_n\omega_n\omega},
\label{eq:modal}
\end{equation}
da cui la condizione di risonanza: quando $\omega \to \omega_{n^\ast}$ il
denominatore del termine $n^\ast$ si riduce a
$2\mathrm{j}\zeta_{n^\ast}\omega_{n^\ast}^2$ e l'ampiezza \`e amplificata di
un fattore $Q = 1/(2\zeta)$ rispetto alla risposta quasi-statica dello stesso
modo ($Q = 100$ per $\zeta = 0{,}5\,\%$). La scelta di $x_F$ e $x_o$ deve
evitare i nodi di $\phi_{n^\ast}$, altrimenti il prodotto
$\phi_{n^\ast}(x_o)\phi_{n^\ast}(x_F)$ annulla la partecipazione del modo:
nello script $x_F = 0{,}30\,L$ e $x_o = 0{,}37\,L$.

\subsection{Spread di risonanza: banda a met\`a potenza}
Dalla \eqref{eq:modal}, l'ampiezza scende di $\SI{3}{dB}$ rispetto al picco
alle frequenze $f_{n^\ast}(1 \pm \zeta)$; la banda a met\`a potenza vale
quindi
\begin{equation}
\Delta f_{-3\,\mathrm{dB}} \simeq 2\,\zeta\, f_{n^\ast}.
\label{eq:band}
\end{equation}
Questo \`e lo \emph{spread} fisicamente rilevante: per
$f_0 = \SI{2}{kHz}$ e $\zeta = 0{,}5\,\%$ si ottiene
$\Delta f \simeq \SI{20}{Hz}$. Se l'eccitazione (o la tensione della campata,
tramite la \eqref{eq:disp}) deriva entro questa banda, la struttura risponde
ancora in condizioni sostanzialmente risonanti; al di fuori, l'ampiezza crolla
rapidamente. \`E anche il criterio con cui interpretare lo sweep numerico
$f \in [0{,}9\,f_0,\;1{,}1\,f_0]$ dello script: il picco \`e stretto
($Q$ elevato) ed \`e la firma inequivocabile della risonanza del modo
$n^\ast$.

\subsection{Transitorio di aggancio alla risonanza}
Integrando la \eqref{eq:pde} nel tempo (Newmark, accelerazione media,
$\gamma = 1/2$, $\beta = 1/4$, incondizionatamente stabile) con eccitazione
esattamente a $f_0$, l'inviluppo della risposta cresce come
\begin{equation}
w_{\mathrm{env}}(t) \propto
\bigl(1 - e^{-\zeta\omega_0 t}\bigr),
\end{equation}
con costante di tempo $\tau = 1/(\zeta\omega_0) \simeq \SI{16}{ms}$: il regime
si raggiunge in $\sim 3\tau \approx \SI{50}{ms}$, cio\`e un centinaio di cicli
a \SI{2}{kHz}. Anche questo \`e un indicatore diagnostico utile: una crescita
lenta dell'inviluppo con saturazione limitata dallo smorzamento \`e tipica
dell'eccitazione \emph{in} risonanza, mentre fuori banda il regime \`e
raggiunto quasi subito con ampiezza modesta e battimenti a
$|f - f_{n^\ast}|$.

\section{Caso numerico di riferimento}
\label{sec:num}
Filo di acciaio armonico: $L = \SI{0.8}{m}$, $d = \SI{2}{mm}$,
$E = \SI{210}{GPa}$, $\rho = \SI{7850}{kg/m^3}$, da cui
$\rho A = \SI{2.47e-2}{kg/m}$ ed $EI = \SI{0.165}{N.m^2}$. Applicando la
\eqref{eq:tuning} con $f_0 = \SI{2}{kHz}$, il modo pi\`u vicino risulta
$n^\ast = 9$ e la tensione di sintonia
\begin{equation*}
T^\ast \simeq \SI{2912}{N}
\qquad (\sigma \simeq \SI{927}{MPa}),
\end{equation*}
ammissibile per filo armonico ad alta resistenza. La
Tabella~\ref{tab:modi} riporta il confronto FEM/analitico ottenuto con lo
script ($N_e = 200$): si noti la non-armonicit\`a dello spettro
($f_2/f_1 > 2$, ecc.) dovuta al termine flessionale.

\begin{table}[ht]
\centering
\caption{Frequenze naturali: FEM (Hermite, $N_e=200$) vs.\ analitico
\eqref{eq:disp}, $T = T^\ast$.}
\label{tab:modi}
\begin{tabular}{ccc}
\toprule
$n$ & $f_n$ FEM [\si{Hz}] & $f_n$ analitica [\si{Hz}]\\
\midrule
1 & 214{,}85 & 214{,}85\\
2 & 430{,}26 & 430{,}26\\
3 & 646{,}79 & 646{,}79\\
4 & 865{,}00 & 865{,}00\\
5 & 1085{,}43 & 1085{,}43\\
6 & 1308{,}63 & 1308{,}63\\
7 & 1535{,}12 & 1535{,}12\\
8 & 1765{,}41 & 1765{,}41\\
\textbf{9} & \textbf{2000{,}00} & \textbf{2000{,}00}\\
10 & 2239{,}38 & 2239{,}38\\
\bottomrule
\end{tabular}
\end{table}

Lo sweep FRF conferma il picco esattamente a \SI{2000}{Hz} con banda a
$-3\,\mathrm{dB}$ di $\approx \SI{20}{Hz}$, in accordo con la
\eqref{eq:band}; l'ampiezza al picco supera di circa due ordini di grandezza
la risposta quasi-statica del modo, coerentemente con $Q = 100$.

\section{Note di validit\`a e possibili estensioni}
\begin{itemize}
\item \textbf{Eulero--Bernoulli vs.\ Timoshenko.} Per il modo 9 la
semilunghezza d'onda \`e $L/9 \approx \SI{89}{mm} \gg d = \SI{2}{mm}$:
taglio ed inerzia rotatoria sono trascurabili e il modello E--B \`e adeguato
fino ben oltre \SI{2}{kHz} per questa geometria.
\item \textbf{Linearit\`a.} Il modello assume $T$ costante; per grandi
ampiezze la tensione dinamica indotta dallo stiramento introduce una
non-linearit\`a di tipo Duffing (irrigidente) che sposta il picco e genera il
caratteristico salto d'isteresi. Il modello lineare resta valido finch\'e
$w_{\max} \ll d$.
\item \textbf{Eccitazione distribuita.} Per forze di volume (es.\ forze
elettromagnetiche) basta sostituire il carico puntuale con la proiezione
$\bm{F} = \int \bm{N}^{\mathsf T} f(x)\,\mathrm{d}x$; la
\eqref{eq:modal} mostra che conta solo la proiezione modale
$\int f(x)\phi_{n^\ast}(x)\,\mathrm{d}x$.
\item \textbf{Incertezza sulla tensione.} Dalla \eqref{eq:disp}, una
variazione relativa $\delta T/T$ sposta $f_{n^\ast}$ di circa
$\tfrac{1}{2}\,\delta T/T$ (nel limite di corda): la condizione perch\'e la
risonanza resti agganciata entro la banda \eqref{eq:band} \`e
$\delta T/T \lesssim 4\zeta$, un vincolo pratico utile in fase di
identificazione sperimentale.
\end{itemize}

\begin{thebibliography}{9}
\bibitem{geradin} M.~G\'eradin, D.~Rixen, \emph{Mechanical Vibrations: Theory
and Application to Structural Dynamics}, 3rd ed., Wiley, 2015.
\bibitem{rao} S.~S.~Rao, \emph{Vibration of Continuous Systems}, 2nd ed.,
Wiley, 2019.
\bibitem{wickert} J.~A.~Wickert, C.~D.~Mote, ``Classical vibration analysis of
axially moving continua,'' \emph{ASME J. Appl. Mech.}, vol.~57, pp.~738--744,
1990.
\bibitem{cook} R.~D.~Cook, D.~S.~Malkus, M.~E.~Plesha, R.~J.~Witt,
\emph{Concepts and Applications of Finite Element Analysis}, 4th ed., Wiley,
2002.
\end{thebibliography}

\end{document}
